\documentclass[11pt]{article}
\usepackage[T1]{fontenc}\usepackage[utf8]{inputenc}
\usepackage[margin=1.05in]{geometry}
\usepackage{amsmath,amssymb,amsthm,mathtools,lmodern,microtype}
\usepackage[authoryear,round]{natbib}\usepackage{xurl}
\usepackage[colorlinks=true,allcolors=blue]{hyperref}
\newcommand{\workbenchurl}{https://github.com/KokunoYumeto/schemes-and-moduli-workbench}
\newcommand{\recordurl}{https://github.com/KokunoYumeto/mathematics-workbenches}
\newcommand{\Fp}{\mathbf F_p}\newcommand{\OO}{\mathcal O}
\DeclareMathOperator{\Spec}{Spec}\DeclareMathOperator{\GL}{GL}
\DeclareMathOperator{\SL}{SL}\DeclareMathOperator{\Sp}{Sp}
\DeclareMathOperator{\End}{End}\DeclareMathOperator{\Hom}{Hom}
\DeclareMathOperator{\Ext}{Ext}\DeclareMathOperator{\Tor}{Tor}
\newtheorem{theorem}{Theorem}[section]
\newtheorem{lemma}[theorem]{Lemma}\newtheorem{proposition}[theorem]{Proposition}
\newtheorem{corollary}[theorem]{Corollary}
\theoremstyle{definition}\newtheorem{definition}[theorem]{Definition}
\theoremstyle{remark}\newtheorem{remark}[theorem]{Remark}
\title{Congruence sections and non-split co-Lie complexes after Godeaux and Serre}
\author{}\date{}
\begin{document}\maketitle
\begin{abstract}
The classical linear congruence-section classifications are due to Sah and Ginosar, and Detinko, Flannery, and Hulpke prove the higher-rank symplectic nonsupplement theorem. We give elementary proofs of the complete section classifications for reduction from $\mathbb Z/p^2$ to $\mathbf F_p$ in general linear, special linear, and symplectic groups. Explicit small-prime obstructions and finite subgroups handle every exceptional case. We then use the Godeaux--Serre construction to realize a nonlifting representation as the nonzero extension class of a co-Lie complex on a smooth projective surface. The identification of the co-Lie complex is due to Bergh and Rydh. Explicit invariant equations and a coherent cohomology injection give full general-linear monodromy in every nonsplit case, while the class in the Grothendieck group of perfect complexes is zero. For primes at least five, rank two is enough. Dimension two is minimal: every such co-Lie complex on a smooth projective curve splits. The mathematics and text were produced and checked by GPT-6 Astra (ChatGPT Pro mode), including separate same-model derivations; no human has checked them.
\end{abstract}

\section{From a representation to a coherent obstruction}

The co-Lie complex of a finite flat group scheme records both invariant differentials and first-order relations at its identity. Even when its two cohomology bundles are isomorphic, the complex can contain an extension between them. Our purpose is to exhibit this extension on a smooth projective surface and to explain why full monodromy by itself does not detect it.

Let $\Gamma$ be a finite group, let $p$ be a prime, and let $V$ be a finite-dimensional $\Fp$-representation of $\Gamma$. Its derived reduction
\[
 K_V=V\otimes^{\mathbf L}_{\mathbb Z}\Fp
       \quad\text{in }D(\Fp[\Gamma])
\]
has cohomology $V$ in degrees $-1$ and $0$. The connecting arrow in its truncation triangle is a class
\begin{equation}\label{eq:kappa}
 \kappa_V\in\Ext^2_{\Fp[\Gamma]}(V,V)
       =H^2(\Gamma,\End_{\Fp}V).
\end{equation}
Thus $K_V$ is isomorphic to $V[1]\oplus V$ if and only if $\kappa_V=0$. Its $\Fp$-module structure comes from the right-hand tensor factor. The two-term resolution $[\mathbb Z\xrightarrow{p}\mathbb Z]$ of that factor computes the underlying $\mathbb Z[\Gamma]$-complex; it does not equip that model with the required derived right-hand $\Fp$-action. The obstruction concerns the object in $D(\Fp[\Gamma])$ with this action. A lift of $V$ means a free $\mathbb Z/p^2$-module whose reduction, including its $\Gamma$-action, is identified with $V$.

If $M$ is an \'etale locally constant sheaf of finite-dimensional $\Fp$-vector spaces on a scheme $S$ of characteristic $p$, let $D(M)$ be its Cartier dual. This is a finite flat group scheme of multiplicative type, \'etale locally isomorphic to a power of $\mu_p$. Write
\[
 E=M\otimes_{\Fp}\OO_S,\qquad
 \ell_{D(M)/S}=e^*L_{D(M)/S},
\]
where $e$ is the identity section and $L$ is the cotangent complex. Bergh and Rydh's formula identifies this complex with
\begin{equation}\label{eq:bergh}
 \ell_{D(M)/S}\simeq M\otimes^{\mathbf L}_{\mathbb Z}\OO_S.
\end{equation}
Its two cohomology sheaves are $E$, and its connecting class lies in
$\Ext^2_S(E,E)=H^2(S,\End E)$.

\begin{theorem}\label{thm:surface}
Suppose that $V$ has no lift to a free $\mathbb Z/p^2$-representation of $\Gamma$. There are an algebraically closed field $k$ of characteristic $p$, a smooth connected projective surface $S/k$, and a connected finite \'etale $\Gamma$-torsor $Y\to S$ such that, for the associated local system $M=Y\times^\Gamma V$, the complex $\ell_{D(M)/S}$ is not isomorphic to $E[1]\oplus E$.

The torsor has surjective monodromy $\pi_1^{\mathrm{et}}(S)\to\Gamma$. In particular, if the representation $V$ is faithful, the monodromy of $M$ is exactly $\Gamma$. Explicit equations for $Y$ over a purely transcendental extension of $\Fp$ are given in Section~\ref{sec:equations}. Nevertheless
\[
 [\ell_{D(M)/S}]=0\quad\text{in }K_0(S).
\]
\end{theorem}

\begin{theorem}[Congruence sections]\label{thm:sections}
Let $p$ be a prime, and let $d,g\geq1$. A section means a group homomorphism whose composite with reduction is the identity. The reduction maps have sections exactly in the following cases:
\[
\begin{array}{c|l}
\text{reduction map}&\text{all split cases}\\ \hline
\GL_d(\mathbb Z/p^2)\longrightarrow\GL_d(\Fp)
 &d=1\text{ or }(d,p)=(2,2),(2,3),(3,2),\\
\SL_d(\mathbb Z/p^2)\longrightarrow\SL_d(\Fp)
 &d=1\text{ or }(d,p)=(2,3),(3,2),\\
\Sp_{2g}(\mathbb Z/p^2)\longrightarrow\Sp_{2g}(\Fp)
 &(g,p)=(1,3).
\end{array}
\]
The symplectic groups preserve the standard nondegenerate alternating form.
\end{theorem}

\begin{corollary}\label{cor:linear}
For the tautological representation of $\GL_d(\Fp)$, Theorem~\ref{thm:surface} applies exactly outside the general-linear split list. Equivalently, it applies for
\[
 p=2,\ d\geq4;\qquad p=3,\ d\geq3;\qquad
 p\geq5,\ d\geq2.
\]
In the split cases the associated co-Lie complex is split on every base. For every $p\geq5$, the surface theorem also applies to the two-dimensional representation of $C_p$ in which a generator acts as $1+E_{12}$. Rank two is minimal for this last assertion: every rank-one representation lifts.
\end{corollary}

\begin{theorem}\label{thm:curve}
For every finite group $\Gamma$ and every $\Fp$-representation $V$, there are a smooth connected projective curve $C$ over an algebraically closed field of characteristic $p$ and a connected finite \'etale $\Gamma$-torsor over $C$. For its associated local system the co-Lie complex is isomorphic to $E[1]\oplus E$, including when $V$ has no $\mathbb Z/p^2$-lift. More generally this splitting holds for every such local system on every smooth projective curve. Dimension two is consequently minimal for the nonvanishing in Theorem~\ref{thm:surface}.
\end{theorem}

These statements distinguish a representation-theoretic obstruction from its image in coherent cohomology. The first obstructs a lift of monodromy; the second obstructs the splitting of a complex of vector bundles. A surjection of fundamental groups does not make the map between these obstruction groups injective. The surface construction enforces injectivity by arranging that the torsor has no coherent cohomology in degree one. The curve construction shows exactly why an unqualified inference from monodromy fails.

\subsection*{Sources, scope, and proof mechanism}

Serre's Proposition~15 in \S20, pp.~50--51 of \citet{Serre}, constructs a smooth complete intersection of any prescribed positive dimension with a free action of a finite group, explicitly crediting Godeaux's method. The geometric existence assertion is therefore classical. We give a particular degree-$|\Gamma|$ norm system and its tangent calculation because they specify equations and work in wild characteristic without averaging. The cohomological survival argument requires only a complete-intersection surface.

Formula~\eqref{eq:bergh} is \citet[Proposition~B.4, pp.~33--34]{BerghRydh}, in their May 2019 version; their Remark~B.6 discusses the representations on the two cohomology groups. The passage from the lifting obstruction to a non-split co-Lie complex uses this formula. The transvection obstruction used in Corollary~\ref{cor:linear} is the argument of \citet[Theorem~6.2, pp.~12--13]{Stasinski}; its calculation works in the general linear group as well as the special linear group. We repeat it and supply all small-prime calculations in Section~\ref{sec:classification}.

The congruence-section lists are classical. \citet[p.~3]{Singla} states the general-linear list and attributes it to Sah and Ginosar. \citet[\S6, p.~12]{Stasinski} attributes the prime-field special-linear classification to \citet[Theorem~7]{Sah74}. \citet[Theorem~2.5, pp.~4--6]{DFH} proves, in particular, that the symplectic congruence kernel in degree at least four has no proper supplement. Their four special-linear exceptions concern proper supplements, which need not be complements. Our proof uses neither their classification nor its finite computations: all nonsplitting cases follow from the identities below, and the positive cases have explicit injectivity certificates. The stated page locators for Stasinski and Detinko--Flannery--Hulpke refer to the specified arXiv versions in the bibliography.

\citet[Theorem~1.2, p.~2, and its proof in \S6, p.~18]{ABM} approximate classifying stacks by smooth projective varieties while preserving Hodge classes in a range. Their construction extends the Godeaux--Serre method to finite group schemes and geometrically reductive groups. Their examples concern nondegeneration of Hochschild and related spectral sequences. Here the finite group is constant, the class has coefficients $\End(V)$, and a degree-two edge map is enough. \citet[Lemma~10.5, pp.~62--63, and Corollary~10.7, pp.~65--66]{Petrov} describes the cotangent complex of the classifying stack of an abelian scheme's $p$-torsion and constructs liftable varieties with a nondegenerate conjugate spectral sequence. His split co-Lie calculation uses an abelian scheme and its multiplication map. Our multiplicative-type group is twisted by a representation that does not lift. These are distinct hypotheses and conclusions.

Recent arithmetic context includes \citet[Theorem~1.1, p.~1]{OpenAI2026}: every continuous irreducible odd representation $G_{\mathbb Q}\to\GL_2(\overline{\mathbb Q}_2)$, unramified outside finitely many primes and de Rham at $2$ with distinct Hodge--Tate weights, is modular up to Tate twist. This characteristic-zero modularity statement is context only; it is not an input to the finite-ring congruence-section or co-Lie arguments.

We make no claim of priority for the synthesis. The results presented here are complete elementary proofs of the congruence-section lists and the precise surface statement above, with explicit equations, examples in every general-linear nonsplit case, and the contrast with curves. The paper makes no assertion about Frobenius degeneration, motives, or a nonzero $K$-theory class. No verified manuscript statement of Grothendieck specific to this co-Lie obstruction is used or quoted.

The geometric proof has three steps. First, a change-of-rings exact sequence identifies $\kappa_V$ with the obstruction to a free lift. Second, norm invariants give a free complete intersection whose $H^1(\OO)$ vanishes. Third, the coherent Cartan--Leray spectral sequence injects the group-cohomological obstruction into $H^2(S,\End E)$. Only the third step distinguishes a surface from a curve.

\section{The lifting obstruction}

This section establishes the algebraic obstruction and the smallest example needed for the theorem.

\begin{lemma}\label{lem:lift}
The class $\kappa_V$ vanishes if and only if $V$ lifts to a free $\mathbb Z/p^2$-representation of $\Gamma$.
\end{lemma}
\begin{proof}
Put $A=\mathbb Z[\Gamma]$ and $B=\Fp[\Gamma]$. Derived extension and restriction of scalars give
\[
 R\Hom_A(V,V)\simeq
 R\Hom_B(B\otimes_A^{\mathbf L}V,V).
\]
The complex $B\otimes_A^{\mathbf L}V$ is $K_V$, because $B=A\otimes_{\mathbb Z}\Fp$ and $A$ is free over $\mathbb Z$. Apply $R\Hom_B(-,V)$ to the triangle
\[
 V[1]\longrightarrow K_V\longrightarrow V
          \xrightarrow{\kappa_V}V[2].
\]
The resulting exact sequence contains
\begin{equation}\label{eq:change}
 \Ext^1_A(V,V)\longrightarrow\Hom_B(V,V)
     \xrightarrow{\partial}\Ext^2_B(V,V),\qquad
 \partial(1_V)=\kappa_V,
\end{equation}
up to the harmless common sign convention for connecting maps.

We identify the first arrow. An element of $\Ext^1_A(V,V)$ is an extension
$0\to V\xrightarrow{i}L\xrightarrow{q}V\to0$ of $A$-modules. Since both outside terms are killed by $p$, multiplication by $p$ on $L$ lands in $i(V)$ and vanishes on $i(V)$. It therefore induces a $B$-linear map $f:V\to V$ characterized by $p\ell=i(f(q(\ell)))$. This is the displayed arrow: applying the resolution $[\mathbb Z\xrightarrow{p}\mathbb Z]$ to the extension identifies its boundary with that formula.

If $f=1_V$, choose a basis $v_1,\ldots,v_d$ of $V$ and lifts $\ell_j$ in $L$. Then $p\ell_j=i(v_j)$. The $\ell_j$ generate $L$ over $\mathbb Z/p^2$: any element first has its image under $q$ removed by a combination of the $\ell_j$, and the remainder is a combination of the $p\ell_j$. If $\sum a_j\ell_j=0$, reduction by $q$ makes each $a_j$ divisible by $p$, and the linear independence of the $i(v_j)$ then makes each $a_j$ zero modulo $p^2$. Thus $L$ is free of rank $d$, with the required action and reduction. Conversely every free lift gives such an extension, using the canonical identification $pL\simeq L/pL$, and its map $f$ is the identity. Exactness of \eqref{eq:change} proves both implications.
\end{proof}

For orientation, choose arbitrary invertible lifts $\widetilde\rho(g)$ of the matrices of the representation. The failure of multiplicativity is
\[
 \widetilde\rho(g)\widetilde\rho(h)
       \widetilde\rho(gh)^{-1}=1+p\,c(g,h).
\]
Associativity says that $c$ is a two-cocycle with conjugation coefficients $\End(V)$. Replacing the lifts changes $c$ by a coboundary. The class vanishes precisely when the lifts can be made multiplicative. This gives the concrete matrix interpretation of the same obstruction; the change-of-rings argument specifies its relation to the complex.

\begin{lemma}\label{lem:transvection}
If $p\geq5$ and $d\geq2$, no matrix in $\GL_d(\mathbb Z/p^2)$ reducing to $U=1+E_{12}$ has order $p$.
\end{lemma}
\begin{proof}
Use the same matrix $U$ over $\mathbb Z/p^2$, and write a lift as $U(1+pX)$. Let $T$ be conjugation by $U$ on $\operatorname{Mat}_d(\Fp)$. The factors $1+pX$ commute with one another modulo $p^2$, so collecting these factors gives
\[
 (U(1+pX))^p=U^p\left(1+p\sum_{i=0}^{p-1}T^{-i}X\right).
\]
Write $N=E_{12}$, so $N^2=0$. Conjugation sends $X$ to
$X+NX-XN-NXN$. Left multiplication by $N$ and right multiplication by $N$ commute and have square zero. Their polynomial expression for $T-1$ consequently has cube zero. In characteristic $p$,
\[
 1+T+\cdots+T^{p-1}=(T-1)^{p-1}=0,
\]
because $p-1\geq4$. The same holds with $T^{-1}$. But $U^p=1+pN\ne1$ modulo $p^2$. This proves the assertion.
\end{proof}

For $\Gamma=C_p$ acting on $\Fp^2$ by $U$, a lift would send its generator to a matrix of order dividing $p$, contradicting the lemma. For the tautological representation of $\GL_d(\Fp)$, restricting a lift to the subgroup generated by $U$ gives the same contradiction. This proves the $p\geq5$ nonlifting assertions in Corollary~\ref{cor:linear}.

For the smallest numerical case $p=5,d=2$, $N=\left(\begin{smallmatrix}0&1\\0&0\end{smallmatrix}\right)$ and
\[
 (1+N)^5=\begin{pmatrix}1&5\\0&1\end{pmatrix}
                       \pmod {25}.
\]
Every possible correction of the lift has zero conjugation norm, so the upper-right entry $5$ cannot be removed. In rank one every character $\Gamma\to\Fp^\times$ lifts by the multiplicative Teichm\"uller section $\Fp^\times\to(\mathbb Z/p^2)^\times$. Lemma~\ref{lem:lift} then gives zero obstruction on every base. The rank-two assertion is therefore sharp.

\section{The exact congruence-section lists}\label{sec:classification}

We prove Theorem~\ref{thm:sections} by explicit matrix calculations. The
negative cases use element orders, one commutation relation, and determinants.
The positive cases are finite subgroups given by matrices, with injectivity
proved directly. No finite-group classification or computer enumeration is
needed.

Throughout this section $R_p=\mathbb Z/p^2$. A bar denotes reduction modulo
$p$, and $E_{ij}$ denotes a matrix unit. For the symplectic group we use
$J_g=\left(\begin{smallmatrix}0&I_g\\-I_g&0\end{smallmatrix}\right)$.

\subsection{Matrix identities and restriction to a summand}

The following formulas keep track of all first-order corrections to a lift.
For integral matrices $U,A$ and a positive integer $m$,
\begin{equation}\label{eq:power-correction}
 (U+pA)^m=U^m+p\sum_{r=0}^{m-1}U^rAU^{m-1-r}\pmod {p^2}.
\end{equation}
Indeed, every term containing two copies of $pA$ vanishes modulo $p^2$.
If $U=I+N$ and $N^s=0$, the sum is
\begin{equation}\label{eq:binomial-correction}
 \sum_{r=0}^{m-1}U^rAU^{m-1-r}
 =\sum_{a,b=0}^{s-1}\binom m{a+b+1}N^aAN^b.
\end{equation}
Here the coefficient identity
$\sum_r\binom ra\binom{m-1-r}b=\binom m{a+b+1}$ counts subsets of an
ordered $m$-element set by their $(a+1)$st element. Coefficients outside their
usual range are zero. We also use
\begin{equation}\label{eq:det-correction}
 \det(U+pA)=\det U\bigl(1+p\operatorname{tr}(U^{-1}A)\bigr)
                       \pmod {p^2},
\end{equation}
which follows by expanding the determinant of $I+pU^{-1}A$.

\begin{lemma}\label{lem:corner}
Let $H\subset\GL_m(\Fp)$ act on a coordinate summand of $\Fp^d$ and
trivially on its complement. A section defined on a group containing this
copy of $H$ gives a lift $H\to\GL_m(R_p)$ by taking the corresponding
matrix corner. If the decomposition is orthogonal for a symplectic form
and the original lift preserves that form, the corner preserves the form
on the summand.
\end{lemma}
\begin{proof}
The lift of $h$ has block form
$\left(\begin{smallmatrix}A(h)&pB(h)\\pC(h)&I+pD(h)\end{smallmatrix}\right)$.
The first corner of a product is $A(h)A(k)$, since the other contribution
contains $p^2$. Its reduction is $h$, so it is invertible. For an orthogonal
decomposition, the first corner of the form equation is
$A(h)^tJ_0A(h)+p^2C(h)^tJ_1C(h)=J_0$.
\end{proof}

A special-linear lift need not have a special-linear corner. We will use
only the general-linear conclusion in that situation.

\subsection{The linear obstructions at the small primes}

Lemma~\ref{lem:transvection} already excludes $\GL_d$ and $\SL_d$ for
$p\geq5$ and $d\geq2$. Its calculation also applies to the symplectic
transvection $I+E_{1,g+1}$. It remains to treat the primes two and three.

\begin{lemma}\label{lem:three-linear}
Neither linear reduction map has a section when $p=3$ and $d\geq3$.
\end{lemma}
\begin{proof}
In $\SL_d(\mathbf F_3)$ the matrices $x=I+E_{12}$ and $y=I+E_{13}$
commute and have order three. Write hypothetical lifts as
$X=I+E_{12}+3A$ and $Y=I+E_{13}+3B$.
Formula~\eqref{eq:power-correction} gives
\[
 X^3=I+3(E_{12}+E_{12}AE_{12})
     =I+3(1+A_{21})E_{12}\pmod9.
\]
Thus $X^3=I$ forces $A_{21}=-1$ in $\mathbf F_3$.
The equation $XY=YX$, however, gives
\[
 AE_{13}-E_{13}A+E_{12}B-BE_{12}=0\quad\text{over }\mathbf F_3.
\]
Its $(2,3)$ entry is $A_{21}=0$, a contradiction. The calculation allows
arbitrary corrections in the full ambient dimension.
\end{proof}

\begin{lemma}\label{lem:two-two}
The reduction $\SL_2(\mathbb Z/4)\to\SL_2(\mathbf F_2)$ has no section.
\end{lemma}
\begin{proof}
Put $U=I+E_{12}$ and $X=U+2A$. The condition $X^2=I$ gives
\[
 E_{12}+E_{12}A+AE_{12}=0\quad\text{over }\mathbf F_2,
 \qquad A_{21}=0,\quad A_{11}+A_{22}=1.
\]
Since $U^{-1}=I-E_{12}$, formula~\eqref{eq:det-correction} now gives
$\det X=1+2(A_{11}+A_{22}-A_{21})=3$ modulo four. Thus the residual
transvection has no involutory lift of determinant one.
\end{proof}

\begin{lemma}\label{lem:four-two}
Put $N=E_{12}+E_{23}+E_{34}$ and $U=I+N$. Every lift $X$ of $\bar U$
over $\mathbb Z/4$ satisfying $X^4=I$ has determinant $-1$.
Consequently neither linear reduction map has a section at $p=2$ in
dimension at least four.
\end{lemma}
\begin{proof}
Write $X=U+2A$. Since $N^4=0$, formulas
\eqref{eq:power-correction}--\eqref{eq:binomial-correction} give
\[
 X^4=I+2N^2+2S,\qquad
 S=N^3A+N^2AN+NAN^2+AN^3\pmod4.
\]
Indeed, the only nonzero coefficient $\binom4{a+b+1}$ modulo two for
$0\leq a,b\leq3$ occurs at $a+b=3$. If $X^4=I$, then $S=N^2$ over
$\mathbf F_2$. Its entries $(1,1),(1,2),(1,3),(1,4)$ give, respectively,
\[
 A_{41}=0,\quad A_{31}+A_{42}=0,\quad
 A_{21}+A_{32}+A_{43}=1,\quad\operatorname{tr}A=0.
\]
Using $U^{-1}=I-N+N^2-N^3$, these equalities yield
\[
 \operatorname{tr}(U^{-1}A)
 =\operatorname{tr}A+(A_{21}+A_{32}+A_{43})
   +(A_{31}+A_{42})+A_{41}=1
 \quad\text{in }\mathbf F_2.
\]
Formula~\eqref{eq:det-correction} proves $\det X=3$.

The group $\GL_4(\mathbf F_2)=\SL_4(\mathbf F_2)$ is perfect.
Elementary matrices generate it by Gaussian elimination; for distinct
$i,j,k$ each one is a commutator
$[I+E_{ik},I+E_{kj}]=I+E_{ij}$.
The determinant of a section would therefore be the trivial homomorphism
to the abelian group $(\mathbb Z/4)^\times$. But $\bar U$ has order four,
so its image would have fourth power one and determinant three. This is
impossible. The group on a four-dimensional summand embeds in every
larger general or special linear group over $\mathbf F_2$.
Lemma~\ref{lem:corner} excludes all those dimensions as well.
\end{proof}

\subsection{The symplectic obstructions}

The cases $p\geq5$ have already been excluded by a transvection. At $p=2$,
restriction to a hyperbolic plane and Lemma~\ref{lem:corner} would give a
section for $\Sp_2=\SL_2$, contradicting Lemma~\ref{lem:two-two}.
The following calculation handles the remaining higher-rank case.

\begin{lemma}\label{lem:sp-three}
The reduction $\Sp_{2g}(\mathbb Z/9)\to\Sp_{2g}(\mathbf F_3)$ has no
section for $g\geq2$.
\end{lemma}
\begin{proof}
First use the symplectic basis $e_1,e_2,f_1,f_2$ and put
\[
 P=I+E_{12}-E_{43},\qquad Q=I+E_{24},\qquad U=PQ=I+N.
\]
Both factors preserve $J_2$: $P=\operatorname{diag}(A,A^{-t})$ with
$A=I+E_{12}$ in dimension two, and $Q$ has a symmetric upper-right block.
Multiplication gives
\[
 N=E_{12}-E_{43}+E_{24}+E_{14},\qquad N^3=-E_{13},\qquad N^4=0.
\]
It follows that $\bar U$ has order nine. For every lift $X=U+3B$,
formula~\eqref{eq:binomial-correction} with $m=9$ shows that its correction
to the ninth power vanishes: all the integers $\binom9k$ for
$1\leq k\leq7$ are divisible by three. Thus
\[
 X^9=U^9=I+9N+36N^2+84N^3=I-3E_{13}\ne I\pmod9.
\]
Even a general-linear lift cannot have the required order. This excludes
$\Sp_4$, and restriction to an orthogonal four-dimensional symplectic
summand excludes every larger $g$.
\end{proof}

\subsection{Rank one and the two-dimensional positive cases}

For $a\in\Fp^\times$ define $s(a)=\widetilde a^{\,p}\bmod p^2$, where
$\widetilde a$ is any integer representative. Changing the representative
by a multiple of $p$ does not change its $p$th power modulo $p^2$.
Consequently $s$ is multiplicative and reduces to $a$; this is the
rank-one section. The group $\SL_1$ is trivial.

Over $\mathbb Z/4$ use
\[
 R=\begin{pmatrix}0&1\\1&0\end{pmatrix},\qquad
 Q=\begin{pmatrix}1&0\\-1&-1\end{pmatrix}.
\]
They satisfy $R^2=Q^2=(RQ)^3=I$. Every word therefore reduces to one of
$(RQ)^j$ and $(RQ)^jR$ for $0\leq j<3$, so the generated group has at
most six elements. Its reduction generates the six-element group
$\GL_2(\mathbf F_2)$: the two distinct involutions have product of order
three. Reduction is thus an isomorphism on this subgroup.

For the two-dimensional case at three, all following matrices are over
$\mathbb Z/9$. Set
\[
 i=\begin{pmatrix}0&1\\-1&0\end{pmatrix},\qquad
 j=\begin{pmatrix}1&4\\4&-1\end{pmatrix},\qquad
 k=ij=\begin{pmatrix}4&-1\\-1&-4\end{pmatrix}.
\]
They obey $i^2=j^2=k^2=-I$ and $ij=-ji$. Their determinants are one,
and the eight elements $\{\pm I,\pm i,\pm j,\pm k\}$ have distinct
reductions. They form a quaternion subgroup $Q_8$.
Define
\[
 h=\frac{-I+i+j+k}{2}=\begin{pmatrix}2&2\\1&6\end{pmatrix}.
\]
For $v=i+j+k$, anticommutation gives $v^2=-3I$, whence
$h^2=(-I-v)/2$ and $h^2+h+I=0$. Thus $h^3=I$, and its reduction has
order three. Multiplication in the quaternion algebra gives
\[
 hih^{-1}=k,\qquad hjh^{-1}=i,\qquad hkh^{-1}=j.
\]
Also $\det h=1$. Hence $H=\langle Q_8,h\rangle$ has at most $24$
elements, while its reduction contains the eight-element quaternion group
and an element of order three normalizing it. The three cosets of that
quaternion group are distinct, so the reduction has $24$ elements. Since
$|\SL_2(\mathbf F_3)|=24$, reduction is an isomorphism on $H$.

To obtain the general-linear section, put
\[
 w=\frac{j-k}{4}=\begin{pmatrix}6&8\\8&3\end{pmatrix}.
\]
The equation $4^2=-2$ modulo nine gives $w^2=I$, and $\det w=-1$.
Direct quaternion multiplication gives
\[
 wiw^{-1}=-i,\qquad wjw^{-1}=-k,\qquad
 wkw^{-1}=-j,\qquad whw^{-1}=h^{-1}.
\]
It follows that $w$ normalizes $H$. The group $\langle H,w\rangle$ has
$48$ elements and reduces onto $\GL_2(\mathbf F_3)$, also of order
$48$, because its image contains $\SL_2(\mathbf F_3)$ and an element
of determinant $-1$. This gives the required section. The identity
$\Sp_2=\SL_2$ gives the positive symplectic case as well.

\subsection{The three-dimensional positive case at two}

Here is a polynomial certificate for the remaining section. Over
$\mathbb Z/4$ set
\[
 S=\begin{pmatrix}0&0&1\\1&0&0\\0&1&0\end{pmatrix},\qquad
 T=\begin{pmatrix}1&1&0\\0&-1&0\\0&2&-1\end{pmatrix},
\]
and consider the homogeneous quartic
\[
 F=x^4+y^4+z^4-x^2y^2-x^2z^2-y^2z^2-xyz(x+y+z).
\]
Both matrices have determinant one. Symmetry of $F$ shows that $S$
preserves it. For $T$ the exact integer identity is
\begin{align*}
 F(x+y,-y,2y-z)-F(x,y,z)
  =4y\bigl(&x^3+x^2y+x^2z+xyz\\
            &+3y^3-7y^2z+6yz^2-2z^3\bigr).
\end{align*}
Thus $T$ also preserves $F$ modulo four.

We claim that the stabilizer of $F$ in $\SL_3(\mathbb Z/4)$ has trivial
reduction kernel. Write an element of that kernel as $I+2A$. Its action
on $F$ is
$F((I+2A)\mathbf x)=F(\mathbf x)+2\,d\bar F_{\mathbf x}(A\mathbf x)$
modulo four. Over $\mathbf F_2$ the three partial derivatives of $\bar F$
are
\[
 yz(y+z),\qquad xz(x+z),\qquad xy(x+y).
\]
In $d\bar F(A\mathbf x)$, the coefficients of
$y^3z,yz^3,x^3z,xz^3,x^3y,xy^3$ force, respectively, all six
off-diagonal entries of $A$ to vanish. The remaining expression is
\[
 xyz\bigl((A_{22}+A_{33})x+(A_{11}+A_{33})y
                         +(A_{11}+A_{22})z\bigr).
\]
It vanishes only if $A=\lambda I$. The determinant condition is then
$\det(I+2\lambda I)=1+2\lambda=1$ modulo four, so $\lambda=0$.
This proves the claim.

Finally, $\bar T=I+E_{12}$. Conjugating by the coordinate cycle $\bar S$
gives $I+E_{23}$ and $I+E_{31}$, and their commutators give the other
three elementary matrices. These generate
$\GL_3(\mathbf F_2)=\SL_3(\mathbf F_2)$ by Gaussian elimination.
Therefore $\langle S,T\rangle$ reduces surjectively, and the stabilizer
claim makes reduction injective. This proves both positive cases.

\begin{remark}\label{rem:seven-lines}
The same injectivity has a short geometric certificate over $\mathbb Z/4$.
The seven free lines generated by
\[
 (1,0,0),\ (0,1,0),\ (0,0,1),\ (1,1,1),\
 (1,2,3),\ (1,3,2),\ (2,1,3)
\]
have distinct nonzero reductions. In this order $S$ permutes them as
$(1\ 2\ 3)(5\ 6\ 7)$ and $T$ as $(2\ 6)(4\ 7)$.
An element in the reduction kernel fixes each line. The coordinate lines
make it diagonal, and the fourth line makes it scalar. Determinant one
then makes it the identity. This gives an alternative to the quartic
calculation, using the same explicit generators.
\end{remark}

\subsection{Assembly and the resulting surface examples}

The preceding arguments prove exactly the lists of
Theorem~\ref{thm:sections}. At $p\geq5$, the transvection excludes every
linear dimension at least two and every positive symplectic rank. At
$p=3$, Lemma~\ref{lem:three-linear} excludes linear dimension at least
three and Lemma~\ref{lem:sp-three} excludes symplectic rank at least two;
the quaternion construction supplies exactly dimension two. At $p=2$,
Lemmas~\ref{lem:two-two} and~\ref{lem:four-two} exclude $\SL_2$ and
all linear dimensions at least four, and the hyperbolic-plane argument
excludes every symplectic group. The explicit sections supply the remaining
linear exceptions, and rank one was handled separately.

For the tautological representation of $\GL_d(\Fp)$, a free
$\mathbb Z/p^2$-lift is precisely a group-theoretic section after choosing
a basis lifting the given basis. Lemma~\ref{lem:lift} therefore makes its
class $\kappa_V$ nonzero exactly outside the general-linear list. Applying
Theorem~\ref{thm:surface} proves the all-prime part of
Corollary~\ref{cor:linear}. In a split case a splitting of the derived
representation descends for every associated local system, so the
co-Lie complex splits on every base.

Preservation of a determinant or an alternating form is an additional
condition on a lift. The special-linear and symplectic nonsplitting lists
alone do not imply nonlifting as a general-linear representation. For
example, the tautological representation of
$\SL_2(\mathbf F_2)=\GL_2(\mathbf F_2)$ has the general-linear lift
constructed above, although no determinant-one lift exists. Our surface
conclusion uses the general-linear classification for precisely this reason.


\section{Explicit free complete intersections}\label{sec:equations}

We now give equations for both the surface and curve torsors. This is the Godeaux--Serre construction with a specified invariant linear system.

Set $m=|\Gamma|$, $D=8m$, and $W_0=\Fp[\Gamma]^{\oplus8}$. Write coordinates of a vector as $(v_{h,j})$ for $h\in\Gamma$ and $1\leq j\leq8$, with $\Gamma$ acting by left translation on $h$. Introduce variables $z_{h,j}$ and define degree-$m$ invariant polynomials $P_\alpha$ by
\begin{equation}\label{eq:norm}
 \prod_{g\in\Gamma}\left(\sum_{h,j}z_{h,j}v_{g^{-1}h,j}\right)
       =\sum_{|\alpha|=m}z^\alpha P_\alpha(v).
\end{equation}
The group permutes the factors, so every coefficient is invariant.

For $r=1$ or $2$, introduce algebraically independent coefficients $c_{a,\alpha}$ for $1\leq a\leq D-1-r$, and put
\begin{align}
 k&=\overline{\Fp(c_{a,\alpha})},\nonumber\\
 Y_r&=V\left(\sum_\alpha c_{a,\alpha}P_\alpha(v) :\
                  1\leq a\leq D-1-r\right)
                   \subset\mathbb P(W_0)_k.\label{eq:Y}
\end{align}
Using independent coefficients means taking the generic point of the full parameter space. No unspecified choice of a rational point in an open subset is needed.

\begin{lemma}\label{lem:free}
The scheme $Y_r$ is a smooth connected complete intersection of dimension $r$, and $\Gamma$ acts freely on it. If $r=2$, then $H^1(Y_2,\OO_{Y_2})=0$.
\end{lemma}
\begin{proof}
First bound the locus with nontrivial projective stabilizer. For a nonidentity element $g$ of order $s$, the permutation of the regular representation consists of cycles of length $s$. On each such cycle every eigenspace has dimension one. This remains true if $p$ divides $s$: the polynomial $T^s-1$ can have repeated roots, but a cyclic permutation has one-dimensional eigenspace for each root. Thus every eigenspace of $g$ on $W_0\otimes\overline{\Fp}$ has dimension $D/s\leq D/2$. Its fixed points in projective space are the union of the projectivizations of these eigenspaces. The nonfree locus $B$ is therefore a finite union of linear subspaces of dimension at most $D/2-1$.

The system spanned by the $P_\alpha$ is base point free. For a nonzero vector $v$, choose a linear form nonzero at all of the finitely many vectors $gv$. Its norm is a member of the system and is nonzero at $v$. Such a form exists over the infinite algebraically closed field.

The system also separates tangent directions at each point of $\mathbb P(W_0)\setminus B$. For a tangent vector represented by $w$, write
\[
 L_g(z)=\sum_{h,j}z_{h,j}v_{g^{-1}h,j},\qquad
 M_g(z)=\sum_{h,j}z_{h,j}w_{g^{-1}h,j}.
\]
If the derivative of the associated projective map is zero, then for some scalar $a$,
\begin{equation}\label{eq:derivative}
 \sum_g M_g\prod_{h\ne g}L_h=a\prod_gL_g.
\end{equation}
Trivial projective stabilizer says that the linear forms $L_g$ are pairwise nonassociate. Reduce \eqref{eq:derivative} modulo the prime ideal $(L_g)$. The other factors remain nonzero in the quotient domain, so $L_g$ divides $M_g$. Taking $g=1$ shows $w$ proportional to $v$. The projective tangent vector was zero.

For completeness, these facts give the required Bertini statement directly by incidence. Put $n=D-1$ and $q=n-r$. At a point of the free locus, the evaluation of the linear system on the value and all $n$ first derivatives is surjective: base point freeness gives the value, and the injective projective differential just proved gives the $n$ cotangent directions. For a tuple of $q$ members, vanishing of all values imposes $q$ independent linear conditions. Failure of rank $q$ for their $q$ differentials has codimension $n-q+1$ in the space of $q$ by $n$ matrices. Thus singular intersection at a specified free point has codimension $n+1$ in the parameter space. Varying the point, of dimension $n$, leaves a subset whose closure is proper. Generic tuples meet transversally throughout the free locus.

At a point of $B$, simultaneous vanishing imposes $q$ independent conditions, since the $q$ members vary independently in a base point free system. As $q=D-1-r>D/2-1\geq\dim B$, the incidence of an intersection with $B$ also has nondense image. The generic tuple avoids $B$. Its intersection is nonempty: $q\leq n$ positive-degree hypersurfaces in projective $n$-space cannot have empty intersection, by the projective dimension theorem (or successive principal-ideal height bounds on the homogeneous cone). Transversality then gives a smooth complete intersection of pure dimension $r$.

Its equations form a regular sequence. The Koszul resolution has terms that are sums of $\OO_{\mathbb P^n}(-jm)$ in position $-j$, $0\leq j\leq q$. For $j>0$ these terms have zero $H^0$, and projective-space cohomology vanishes in degrees strictly between $0$ and $n$. The only additional possible contribution to the hypercohomology of the resolution has total degree at least $n-q=r$. Hence $H^0(Y_r,\OO_{Y_r})=k$, and $H^i(Y_r,\OO_{Y_r})=0$ for $0<i<r$. The first equality proves connectedness. For $r=2$ the second gives the claimed $H^1$-vanishing. Avoidance of $B$ proves freeness.
\end{proof}

The quotient $S_r=Y_r/\Gamma$ exists as a projective scheme: a sufficiently high power of the ample equivariant bundle $\OO(1)$ descends, and invariant affine charts give the finite quotient. Because the constant finite group acts freely, $Y_r\to S_r$ is a finite \'etale $\Gamma$-torsor, including when $p$ divides $|\Gamma|$. Smoothness descends along this \'etale cover. Since $Y_r$ is connected, the monodromy homomorphism onto $\Gamma$ is surjective: a proper subgroup image would decompose the torsor into its distinct cosets.

\section{Survival of the co-Lie class}

This section identifies the complex and uses coherent cohomology to distinguish the surface from the curve.

For the torsor $Y_r\to S_r$, let $M=Y_r\times^\Gamma V$ and $E=M\otimes_{\Fp}\OO_{S_r}$. Formula~\eqref{eq:bergh} holds \'etale locally by the diagonalizable-group calculation of Bergh and Rydh, and it is functorial in the character module. It therefore descends along the torsor with precisely the $\Gamma$-action on $V$. In particular, its truncation triangle is
\begin{equation}\label{eq:triangle}
 E[1]\longrightarrow\ell_{D(M)/S_r}\longrightarrow E
              \xrightarrow{\xi}E[2],
\end{equation}
where $\xi$ is the image of $\kappa_V$ under group-cohomology scalar extension and coherent descent. The two-term complex is perfect: locally $D(M)=\mu_p^{\dim V}$, and the character lattice presentation gives two finite free terms. Its two cohomology sheaves are the vector bundle $E$.

\begin{proposition}\label{prop:edge}
If $H^0(Y,\OO_Y)=k$ and $H^1(Y,\OO_Y)=0$ for a finite \'etale $\Gamma$-torsor $Y\to S$, then for every finite-dimensional $k$-representation $T$ the natural map
\[
 H^2(\Gamma,T)\longrightarrow H^2(S,Y\times^\Gamma T)
\]
is injective.
\end{proposition}
\begin{proof}
Use the \v Cech nerve of the torsor and an injective or affine \v Cech resolution for coherent cohomology on $Y$. The resulting double complex has spectral sequence
\[
 E_2^{a,b}=H^a\bigl(\Gamma,H^b(Y,\OO_Y)\otimes_k T\bigr)
          \Longrightarrow H^{a+b}(S,Y\times^\Gamma T).
\]
This construction uses group cochains, not an exactness assertion for invariants, so it is valid in wild characteristic. The term $E_2^{2,0}=H^2(\Gamma,T)$ has no outgoing differential, because every possible target has negative second index. The only possible incoming differential is $d_2:E_2^{0,1}\to E_2^{2,0}$: for $d_j$ with $j\geq3$ the source has negative first index. The assumed $H^1$-vanishing kills $E_2^{0,1}$. Thus $E_\infty^{2,0}=E_2^{2,0}$. It is the last nonzero filtered piece in total degree two, hence a subspace of $H^2(S,Y\times^\Gamma T)$. This is the stated edge map.
\end{proof}

Apply the proposition with $r=2$ and $T=\End_k(V\otimes_{\Fp}k)$. Group cohomology commutes with extension of the ground field here: the inhomogeneous group-cochain complex in every degree is a finite direct sum of the coefficient vector space, and the field extension is flat. Consequently a nonzero $\kappa_V$ remains nonzero over $k$. Lemmas~\ref{lem:lift} and~\ref{lem:free} and Proposition~\ref{prop:edge} give $\xi\ne0$. The triangle~\eqref{eq:triangle} therefore does not split. This proves the existence and monodromy assertions in Theorem~\ref{thm:surface}. In $K_0(S)$ the same triangle gives
\[
 [\ell_{D(M)/S}]=[E[1]]+[E]=-[E]+[E]=0.
\]
This equality of classes does not split a triangle. It proves the last assertion of that theorem and shows why the obstruction cannot be recovered from this Grothendieck-group class alone.

For $r=1$, Lemma~\ref{lem:free} gives the required connected torsor over a smooth projective curve. For any vector bundle $E$ on a curve,
\[
 \Ext^2_C(E,E)=H^2(C,\End E)=0.
\]
Here the first equality follows because $E$ is locally free, and the second is coherent cohomological dimension one. Thus $\xi=0$, the triangle splits, and Theorem~\ref{thm:curve} follows. The splitting need not be canonical. The same vanishing on dimension-zero smooth projective schemes proves the asserted minimum dimension.

\begin{remark}\label{rem:as}
The loss of information has another exact expression. For any such local system, its associated endomorphism bundle has a Frobenius-semilinear endomorphism $F$, and the \'etale Artin--Schreier sequence is
\[
 0\longrightarrow\End_{\Fp}(M)\longrightarrow\End(E)
       \xrightarrow{F-1}\End(E)\longrightarrow0.
\]
It is exact because \'etale locally the local system is constant and the scalar Artin--Schreier equation is \'etale locally solvable. Its long exact sequence identifies the kernel of
\[
 H^2_{\mathrm{et}}(S,\End M)\longrightarrow H^2(S,\End E)
\]
with the image of the connecting map from
$H^1(S,\End E)/(F-1)H^1(S,\End E)$. Quasi-coherent \'etale and Zariski cohomology agree in this formula. A full monodromy hypothesis gives no reason for that connecting image to vanish.
\end{remark}

\section{Scope of the remaining questions}

The realization theorem applies to every representation with nonzero lifting obstruction. For the tautological general-linear representations, Theorem~\ref{thm:sections} resolves this representation-theoretic test for every prime and dimension. Classifying lifts of arbitrary representations remains a separate problem. Moreover, a nonzero group-cohomological obstruction need not survive on every surface with full monodromy. Our construction supplies a surface on which it survives; it makes no universal nonvanishing assertion. The curve theorem shows why preservation of the degree-two class matters. On a specified surface the precise remaining test is whether the image of $\kappa_V$ in $H^2(S,\End E)$ vanishes.

\section*{Note on method and records}

The mathematics and text were produced by GPT-6 Astra (ChatGPT Pro mode). The surface, curve, and lifting-obstruction arguments received separate same-model derivation and complete comparison. The three congruence-section lists also received two separate derivations, followed by full comparison. The exceptional matrices were shared as statements before the second derivation; their group relations and injectivity certificates were derived separately. Exact arithmetic checks use polynomial identities and a second calculation with seven lifted lines and quaternion normal forms. The written proofs do not rely on enumerated group orders. The effort setting is not exposed. No human has checked the mathematics. No claim of novelty is made for the classical classifications or for these proof methods. The text is dedicated to the public domain under CC0~1.0. Sources and verification records belong to \url{\workbenchurl}; their permanent records address is \url{\recordurl}.

\begin{thebibliography}{99}
\bibitem[Antieau et~al.(2021)]{ABM}
Antieau, B., Bhatt, B., \& Mathew, A. (2021). Counterexamples to Hochschild--Kostant--Rosenberg in characteristic $p$. \emph{Forum of Mathematics, Sigma, 9}, e49, 1--26. \url{https://doi.org/10.1017/fms.2021.20}. Locators in the text refer to the February 2021 author manuscript, \url{https://arxiv.org/abs/1909.11437}.
\bibitem[Bergh and Rydh(2019)]{BerghRydh}
Bergh, D., \& Rydh, D. (2019). \emph{Functorial destackification and weak factorization of orbifolds}. May 2, 2019 version. \url{https://arxiv.org/abs/1905.00872}.
\bibitem[Detinko et~al.(2018)]{DFH}
Detinko, A., Flannery, D. L., \& Hulpke, A. (2018). Zariski density and computing in arithmetic groups. \emph{Mathematics of Computation, 87}(310), 967--986. \url{https://doi.org/10.1090/mcom/3236}. Consulted version and locators: arXiv:1611.05921v4, November 4, 2022, \url{https://arxiv.org/abs/1611.05921v4}.
\bibitem[OpenAI(2026)]{OpenAI2026}
OpenAI. (2026, October 6). \emph{Fontaine--Mazur modularity at the prime 2} [Preprint, commit \texttt{4eaee8754825}]. \url{https://github.com/openai/math/tree/4eaee87548250b2682affda7d05d947c05554bae/preprints/Fontaine-Mazur-modularity-at-the-prime-2-October-6-2026}.
\bibitem[Petrov(2025)]{Petrov}
Petrov, A. (2025). \emph{Non-decomposability of the de Rham complex and non-semisimplicity of the Sen operator}. Version 2, October 12, 2025. \url{https://arxiv.org/abs/2302.11389}.
\bibitem[Sah(1974)]{Sah74}
Sah, C.-H. (1974). Cohomology of split group extensions. \emph{Journal of Algebra, 29}(2), 255--302. \url{https://doi.org/10.1016/0021-8693(74)90099-4}. The classification attribution is recorded in Stasinski's Section~6.
\bibitem[Serre(1958)]{Serre}
Serre, J.-P. (1958). Sur la topologie des vari\'et\'es alg\'ebriques en caract\'eristique $p$. In \emph{Symposium internacional de topolog\'ia algebraica} (pp.~24--53). Universidad Nacional Aut\'onoma de M\'exico and UNESCO. \url{https://webhomes.maths.ed.ac.uk/~v1ranick/papers/mexico.pdf}.
\bibitem[Singla(2010)]{Singla}
Singla, P. (2010). \emph{Representations and conjugacy classes of general linear groups over principal ideal local rings of length two} [Doctoral dissertation, Homi Bhabha National Institute]. Institute of Mathematical Sciences repository. \url{https://www.imsc.res.in/xmlui/handle/123456789/198}.
\bibitem[Stasinski(2021)]{Stasinski}
Stasinski, A. (2021). Representations of $\mathrm{SL}_n$ over finite local rings of length two. \emph{Journal of Algebra, 566}, 119--135. \url{https://doi.org/10.1016/j.jalgebra.2020.08.036}. Consulted version and locators: arXiv:1907.09878v2, \url{https://arxiv.org/abs/1907.09878v2}.
\end{thebibliography}
\end{document}
